% Options for packages loaded elsewhere
\PassOptionsToPackage{unicode}{hyperref}
\PassOptionsToPackage{hyphens}{url}
\PassOptionsToPackage{dvipsnames,svgnames,x11names}{xcolor}
\documentclass[
  10pt,
  fontset=fandol]{ctexart}
\usepackage{xcolor}
\usepackage[margin=16mm]{geometry}
\usepackage{amsmath,amssymb}
\setcounter{secnumdepth}{-\maxdimen} % remove section numbering
\usepackage{iftex}
\ifPDFTeX
  \usepackage[T1]{fontenc}
  \usepackage[utf8]{inputenc}
  \usepackage{textcomp} % provide euro and other symbols
\else % if luatex or xetex
  \usepackage{unicode-math} % this also loads fontspec
  \defaultfontfeatures{Scale=MatchLowercase}
  \defaultfontfeatures[\rmfamily]{Ligatures=TeX,Scale=1}
\fi
\usepackage{lmodern}
\ifPDFTeX\else
  % xetex/luatex font selection
\fi
% Use upquote if available, for straight quotes in verbatim environments
\IfFileExists{upquote.sty}{\usepackage{upquote}}{}
\IfFileExists{microtype.sty}{% use microtype if available
  \usepackage[]{microtype}
  \UseMicrotypeSet[protrusion]{basicmath} % disable protrusion for tt fonts
}{}
\makeatletter
\@ifundefined{KOMAClassName}{% if non-KOMA class
  \IfFileExists{parskip.sty}{%
    \usepackage{parskip}
  }{% else
    \setlength{\parindent}{0pt}
    \setlength{\parskip}{6pt plus 2pt minus 1pt}}
}{% if KOMA class
  \KOMAoptions{parskip=half}}
\makeatother
\setlength{\emergencystretch}{3em} % prevent overfull lines
\providecommand{\tightlist}{%
  \setlength{\itemsep}{0pt}\setlength{\parskip}{0pt}}
\usepackage{amsmath,amssymb,mathtools,needspace}
\xeCJKDeclareCharClass{CJK}{"2032,"0400->"04FF,"2160->"216B,"2460->"2473,"25A0,"25C7,"25CB,"25CF}
\IfFontExistsTF{Arial Unicode MS}{\xeCJKsetup{AutoFallBack=true}\setCJKfallbackfamilyfont{\CJKrmdefault}{Arial Unicode MS}}{}
\setcounter{secnumdepth}{0}
\setlength{\emergencystretch}{2em}
\allowdisplaybreaks
\usepackage{bookmark}
\IfFileExists{xurl.sty}{\usepackage{xurl}}{} % add URL line breaks if available
\urlstyle{same}
\hypersetup{
  pdftitle={误差与误差限},
  colorlinks=true,
  linkcolor={Maroon},
  filecolor={Maroon},
  citecolor={Blue},
  urlcolor={Blue},
  pdfcreator={LaTeX via pandoc}}

\title{误差与误差限}
\author{}
\date{}

\begin{document}
\maketitle

\paragraph{1.3.1 误差与误差限}\label{ux8befux5deeux4e0eux8befux5deeux9650}

\textbf{定义1.1} 设 \(x\) 为准确值，\(x^*\) 为 \(x\) 的一个近似值，称 \(e^* = x^* - x\) 为近似值的绝

\begin{quote}
校注（agent 补充）：本页定义1.1末尾``绝''字与下一页``对误差''相接。
\end{quote}

\begin{quote}
校注（agent 补充，CH01-E001）：原书写``则 \(\tilde{I}_n\) 就是 \(E_0\) 的 \(n!\) 倍误差''，此处照录。由上式 \(E_n=(-1)^n n!E_0\) 可知，应理解为近似值 \(\tilde{I}_n\) 的误差绝对值为初值误差绝对值的 \(n!\) 倍；疑为原书表述不严谨，不是 OCR 错误。
\end{quote}

\href{../../source-images/pdf-018.jpeg}{查看原页}

对误差, 简称误差.

注意, 这样定义的误差 \(e^*\) 可正可负, 当绝对误差为正时近似值偏大, 叫做强近似值; 当绝对误差为负时近似值偏小, 叫做弱近似值.

通常, 我们不能算出准确值 \(x\), 也不能算出误差 \(e^*\) 的准确值, 只能根据测量工具或计算情况估计出误差的绝对值不超过某正数 \(\varepsilon^*\), 也就是误差绝对值的一个上界. \(\varepsilon^*\) 叫做近似值的误差限, 它总是正数. 例如, 用毫米刻度的米尺测量一长度 \(x\) (单位: mm, 下同), 读出和该长度接近的刻度 \(x^*\), \(x^*\) 是 \(x\) 的近似值, 它的误差限是 0.5, 于是 \(|x^* - x| \leqslant 0.5\); 如读出的长度为 765, 则有 \(|765 - x| \leqslant 0.5\). 从该不等式仍不知道准确的 \(x\) 是多少, 但知道 \(764.5 \leqslant x \leqslant 765.5\), 说明 \(x\) 在区间 \([764.5, 765.5]\) 上.

对于一般情形, \(|x^* - x| \leqslant \varepsilon^*\), 即 \(x^* - \varepsilon^* \leqslant x \leqslant x^* + \varepsilon^*\), 这个不等式有时也表示为

\[
x = x^* \pm \varepsilon^*.
\]

\subsection{本节覆盖原页的校注记录}\label{ux672cux8282ux8986ux76d6ux539fux9875ux7684ux6821ux6ce8ux8bb0ux5f55}

下列记录按原页关联，含同页相邻小节的校注；计算时同时读取上文校注和完整原页。

\begin{itemize}
\tightlist
\item
  \texttt{CH01-E001}：原书 PDF 页 17
\end{itemize}

\end{document}
